\documentclass{article}\usepackage[T1]{fontenc}\begin{document}
For the common-size, all-length telescope over an ordinary base with the existing split endpoints and a fixed ordinary Cat-valued presheaf, evaluationUnitDiagram and evaluationCompDiagram prove the identity and composition equations for the actual SpanPull.All.evaluationBase natural isomorphisms. The source constraints are those of the chosen weighted cleavage of the single coproduct span. The target constraints come from the existing tagged iterate presheaf, not a newly defined action. The composition diagram pulls back along the second arrow and then the first. All lengths, including zero, occur in the same functor. This closes the two base diagrams at fixed payload; it does not prove mixed coherence with changes of payload for this all-length single-span comparison, bundled higher coherence, flattening or strict algebra recovery.

For the source route, read the manuscript at 458ff36a, lines 929--953 (eq:telescope), then SpanPull.All.evaluationBase and its two operands. \#199 proved the forward square laws before replacing the generic component-pull family with the concrete length-indexed family. \#200 proved the source edge has the same unit and compositor after its specified equality transport. In MateLaws, adjustUnit and adjustComp move the old forward diagrams across this correction; All.comparisonUnit and All.comparisonComp concern the actual corrected forward square. In mate\_fac an arrow into an inverse functor is tested by mapping forward and composing with the chosen counit. This test is injective by faithfulness and cancellation. Expanding the actual inverseSquare and using the equivalence identities computes its test; the result retains the forward cell and the image of the starting counit. For the composite, naturality of the first forward cell at the second mate is essential. The ordinary functoriality of pull alone does not identify those paths. All.mateUnit and All.mateComp then paste with \#200 at the same length-indexed constraints. The final natural-transformation equations use exactly the old evaluationBase. If the choice of mate, the edge equality, or either action changes, recheck mate\_fac and both bridge conditions before reusing the paste.
\end{document}
